\documentclass[10pt,a4paper,twoside]{article}
\usepackage[utf8]{inputenc}
\usepackage[vietnamese]{babel}
\usepackage{amsmath,amssymb,amsfonts,amsthm}
\usepackage{booktabs}
\usepackage{multirow}
\usepackage{graphicx}
\usepackage{geometry}
\geometry{a4paper, margin=16mm, top=18mm, bottom=20mm, headheight=14pt}
\usepackage{hyperref}
\usepackage{caption}
\usepackage{subcaption}
\usepackage{float}
\usepackage{microtype}
\usepackage{fancyhdr}
\usepackage{enumitem}
\usepackage{algorithm}
\usepackage{algpseudocode}
\usepackage{xcolor}

\hypersetup{colorlinks=true, linkcolor=blue!80!black, citecolor=blue!80!black, urlcolor=blue!80!black}

\pagestyle{fancy}
\fancyhf{}
\fancyhead[CE]{\small\textsc{Nhóm Nghiên Cứu Machine Learning \& Khai Phá Dữ Liệu}}
\fancyhead[CO]{\small\textsc{GSI-MLC-PA v6.3.2: One-Step Normalized Mean-Field Coupling}}
\fancyfoot[C]{\small\thepage}
\renewcommand{\headrulewidth}{0.4pt}
\renewcommand{\footrulewidth}{0pt}

\fancypagestyle{firstpage}{
    \fancyhf{}
    \fancyhead[L]{\small\textbf{Báo Cáo Nghiên Cứu Khoa Học Máy Tính (Antigravity Research)}}
    \fancyhead[R]{\small\textit{Thực nghiệm độc lập; Xuất bản 10/2026}}
    \fancyfoot[L]{\footnotesize\copyright 2026 Nhóm Nghiên Cứu Machine Learning. Độc quyền thực nghiệm.}
    \fancyfoot[R]{\small 1}
    \renewcommand{\headrulewidth}{0.4pt}
    \renewcommand{\footrulewidth}{0pt}
}

\begin{document}
\thispagestyle{firstpage}

\begin{center}
    {\LARGE\bfseries Suy Diễn Biến Phân Trường Trung Bình Chuẩn Hóa Một Bước Trên Ma Trận Tương Quan Phần Dư\\ Trong Phân Loại Đa Nhãn Có Chọn Lọc:\\[0.3em] Đánh Giá Thực Nghiệm Toàn Diện Mô Hình GSI-MLC-PA v6.3.2}\\[0.6em]
    {\large\textit{One-Step Normalized Mean-Field Variational Inference on Residual Error Correlation for Selective\\ Multi-Label Classification: Empirical Study on GSI-MLC-PA v6.3.2 across Multiple Base Classifiers}}\\[0.9em]
    {\textbf{Nhóm Nghiên Cứu Machine Learning \& Khai Phá Dữ Liệu}}\\[0.2em]
    {\texttt{ml.research@lab.edu.vn}}\\[0.2em]
    {\small Phòng Thí Nghiệm Trí Tuệ Nhân Tạo Nâng Cao \& Khai Phá Dữ Liệu}\\[0.1em]
    {\small Khoa Khoa Học Máy Tính, Trường Đại học}\\[0.8em]
\end{center}

\begin{center}\textbf{Tóm tắt (Abstract)}\end{center}
Trong bài toán phân loại đa nhãn có chọn lọc (Selective Multi-Label Classification), việc mô hình hóa các nhãn phụ thuộc ($DL$) thường đòi hỏi huấn luyện các bộ phân loại điều kiện thứ cấp $g_l(x, \hat{P}_{\text{cha}})$. Tuy nhiên, phương pháp này bộc lộ ba hạn chế kỹ thuật: gia tăng chi phí tính toán, nguy cơ lan truyền sai số chuỗi khi số lượng nhãn phụ thuộc lớn ($|DL| \ge 10$), và thiếu ràng buộc đại số nhất quán giữa các nhãn tương quan. Để khắc phục triệt để các hạn chế này, chúng tôi đề xuất kiến trúc \textbf{GSI-MLC-PA v6.3.2} với phương pháp \textbf{Suy diễn biến phân Trường trung bình Chuẩn hóa một bước (One-Step Normalized Mean-Field - OS-NMF)}. Phương pháp khai thác trực tiếp ma trận tương quan Pearson có dấu của phần dư sai số ngoại mẫu ($C^{\text{res}}$), đồng thời tích hợp ba cơ chế kiểm soát chặt chẽ: (1) \textit{Chuẩn hóa bậc kết nối} $\max(1.0, \sum_m |C^{\text{res}}_{jm}|)$ nhằm ngăn chặn hiện tượng bão hòa xác suất và đếm lặp phụ thuộc; (2) \textit{Căn giữa spin đối xứng} $2P - 1.0$ bảo toàn cân bằng kỳ vọng thống kê dưới điều kiện mất cân bằng nhãn; và (3) \textit{Kẹp biên độ logit} $[-z_{\max}, +z_{\max}]$ (với $z_{\max} = 0.50$) bảo đảm tương thích trọn vẹn với chốt chặn Precision Guard ($\tau_1 \ge 0.50$) và hiệu chuẩn Balanced-Root Platt. Thực nghiệm đối chuẩn 5-Fold Cross-Validation trên \textbf{10 tập dữ liệu benchmark quốc tế} với \textbf{3 bộ học cơ sở} (Logistic Regression, Calibrated Linear SVM, MLP) --- tương ứng 30 cấu hình kiểm định độc lập --- đối đầu trực tiếp với BR, CC, MLC-PA, GSI v6.2.1 và GSI v6.3.1 ghi nhận: (i) \textbf{Selective Macro-$F_1$} đạt \textbf{0.5317} (vượt trội các mô hình chuẩn); (ii) \textbf{Subset 0/1 Accuracy} tăng vọt lên \textbf{0.3425}; (iii) \textbf{Selective Hamming Loss} giảm sâu về \textbf{0.1515}; (iv) Duy trì độ bao phủ an toàn vững chắc \textbf{73.3\%}.

\vspace{0.5em}
\noindent\textbf{Từ khóa:} Phân loại đa nhãn (MLC), Dự đoán có chọn lọc (Selective Classification), Mean-Field Variational Inference, Tương quan phần dư (PCC), Subset Accuracy, Precision Guard, Base Learners.

\section{Giới Thiệu (Introduction)}
Trong phân loại đa nhãn (Multi-Label Classification - MLC), bài toán dự đoán có chọn lọc (Selective Prediction) cho phép mô hình từ chối đưa ra phán đoán trên các nhãn có độ bất định cao nhằm bảo vệ độ chính xác trên tập các mẫu được chấp nhận [4, 6]. Dòng mô hình GSI-MLC-PA sử dụng chiến lược bóc tách phân tầng kiểm định chéo (Stratified CV Peeling) để chia nhãn thành hai nhóm: nhóm nhãn độc lập ($IL$) và nhóm nhãn phụ thuộc dư thừa ($DL$) [5].

Ở phiên bản \textbf{v6.3.1}, việc đưa vào cơ chế Hiệu chuẩn Căn bậc hai (Balanced-Root Platt) và Chốt chặn Precision Guard ($\tau_1 \ge 0.50$) đã dập tắt hiện tượng lạm phát xác suất trên nhãn hiếm. Tuy nhiên, việc mô hình hóa các nhãn phụ thuộc $DL$ vẫn phụ thuộc vào việc huấn luyện các bộ phân loại điều kiện $g_l(x, \hat{P}_{\text{cha}})$. Phương pháp này bộc lộ những hạn chế rõ rệt về tài nguyên tính toán và tính ổn định khi số lượng nhãn phụ thuộc lớn ($|DL| \ge 10$). Bài báo này giới thiệu phiên bản hoàn thiện \textbf{GSI-MLC-PA v6.3.2}, thay thế toàn bộ các mô hình điều kiện $g_l$ bằng suy diễn biến phân trường trung bình chuẩn hóa một bước (OS-NMF) dựa trên ma trận tương quan Pearson của phần dư sai số, đồng thời đánh giá toàn diện trên 3 họ phân loại cơ sở: Logistic Regression, Linear SVM và Multilayer Perceptron (MLP).

\section{Hạn Chế Của Mô Hình Phân Loại Điều Kiện Thứ Cấp ở v6.3.1}
Trong v6.3.1, với mỗi nhãn $l \in DL$, thuật toán tìm tập nhãn cha $DL_{\text{temp}}[l]$ có tương quan phần dư vượt ngưỡng $\theta_{\text{corr}} = 0.25$, sau đó ghép vector xác suất ngoài mẫu của các nhãn cha vào không gian đặc trưng để huấn luyện bộ phân loại nhị phân thứ cấp $g_l$. Mô hình này tồn tại 3 nhược điểm cơ bản:
\begin{itemize}[leftmargin=*]
    \item \textbf{Gánh nặng tính toán:} Phải huấn luyện thêm $|DL|$ bộ phân loại trên không gian đặc trưng mở rộng $[X, P_{DL_{\text{temp}}}]$. Trên các tập dữ liệu có $|DL| = 14$ như \texttt{humanpseaac}, chi phí huấn luyện tăng gấp đôi.
    \item \textbf{Lan truyền sai số chuỗi (Error Propagation):} Nếu các nhãn cha bị dự đoán sai ở bước cơ sở, vector xác suất sai lệch sẽ trở thành đặc trưng đầu vào nhiễu cho $g_l$, gây hiện tượng sai số dây chuyền sang các nhãn con.
    \item \textbf{Mâu thuẫn dự đoán giữa các nhãn liên đới:} Mô hình $g_l$ huấn luyện độc lập cho từng nhãn con, không có cơ chế cân bằng đại số hai chiều, dẫn đến việc các nhãn có tương quan đồng biến thiên mạnh vẫn có thể bị dự đoán trái ngược nhau.
\end{itemize}

\section{Kiến Trúc Hoàn Thiện: GSI-MLC-PA v6.3.2}
Thay vì huấn luyện các mô hình điều kiện phức tạp, GSI-MLC-PA v6.3.2 áp dụng cơ chế \textbf{Suy diễn biến phân Trường trung bình Chuẩn hóa một bước (One-Step Normalized Mean-Field - OS-NMF)} cập nhật trực tiếp logit của các nhãn trong $DL$.

\subsection{Tính Toán Ma Trận Tương Quan Phần Dư Có Dấu (Signed Residual PCC)}
Gọi $P^{\text{OOF}}_{DL} \in [0, 1]^{N \times |DL|}$ là ma trận xác suất ngoài mẫu từ mô hình BR cơ sở trên không gian ngữ cảnh $X_{\text{context}}$. Ma trận phần dư sai số liên tục được định nghĩa:
\begin{equation}
R_{i, j} = Y_{i, j} - P^{\text{OOF}}_{i, j}, \quad \forall j \in DL.
\end{equation}
Ma trận tương quan Pearson có dấu $C^{\text{res}} \in [-1, 1]^{|DL| \times |DL|}$ giữa các cột phần dư được tính toán:
\begin{equation}
C^{\text{res}}_{j, k} = \frac{\sum_{i=1}^N (R_{i, j} - \bar{R}_j)(R_{i, k} - \bar{R}_k)}{\sqrt{\sum_{i=1}^N (R_{i, j} - \bar{R}_j)^2} \sqrt{\sum_{i=1}^N (R_{i, k} - \bar{R}_k)^2}}, \quad \forall j \neq k.
\end{equation}
Các phần tử có $|C^{\text{res}}_{j, k}| < \theta_{\text{corr}} = 0.25$ được triệt tiêu về $0$.

\subsection{Chuẩn Hóa Bậc Hàng và Căn Giữa Spin Đối Xứng}
Để ngăn ngừa nguy cơ bão hòa xác suất khi nhãn có nhiều liên kết, ma trận trọng số tương quan được chuẩn hóa theo tổng bậc hàng:
\begin{equation}
W_{j, k} = \frac{C^{\text{res}}_{j, k}}{\max\left(1.0, \, \sum_{m \neq j} |C^{\text{res}}_{j, m}|\right)}.
\label{eq:degree_norm}
\end{equation}
Đồng thời, xác suất của các nhãn đối tác được căn giữa sang biến spin Ising đối xứng $s_k = 2 p_k - 1.0 \in [-1, 1]$. Độ dịch chuyển logit thô của nhãn $j$ được tính toán thông qua tương tác trường trung bình:
\begin{equation}
\Delta z_j = \alpha \sum_{k \neq j} W_{j, k} \cdot (2 p_k - 1.0),
\end{equation}
trong đó $\alpha = 0.25$ là hệ số ghép nối trường trung bình.

\subsection{Kẹp Biên Độ Logit và Tinh Chỉnh Xác Suất}
Nhằm bảo đảm độ dịch chuyển logit không phá vỡ tính đúng đắn của phép hiệu chuẩn Platt và chốt chặn Precision Guard, độ dịch chuyển được kẹp chặt trong khoảng $[-z_{\max}, +z_{\max}]$ với $z_{\max} = 0.50$:
\begin{equation}
\Delta z_j^{(\text{clipped})} = \text{clip}\left(\Delta z_j, \, -z_{\max}, \, +z_{\max}\right).
\end{equation}
Xác suất tinh chỉnh cuối cùng của nhãn phụ thuộc $j \in DL$ được xác định qua hàm Sigmoid:
\begin{equation}
p_j^{(\text{refined})} = \sigma\left(\text{logit}(p_j) + \Delta z_j^{(\text{clipped})}\right).
\label{eq:refined_prob}
\end{equation}

\begin{algorithm}[H]
\caption{Quy trình Huấn luyện và Suy diễn Chọn lọc Toàn diện GSI-MLC-PA v6.3.2}
\label{alg:v632}
\begin{algorithmic}[1]
\Require Tập huấn luyện $(X, Y) \in \mathbb{R}^{N \times d} \times \{0, 1\}^{N \times K}$, tập kiểm tra $X^* \in \mathbb{R}^{N^* \times d}$, danh sách ngưỡng bóc tách $\mathcal{T} = [0.75, 0.70, 0.65]$, ngưỡng tương quan phần dư $\theta_{\text{corr}} = 0.25$, chi phí từ chối $c=0.30$, độ phủ sàn $\gamma_{\min}=0.70$, hệ số ghép nối $\alpha=0.25$, biên độ kẹp $z_{\max}=0.50$.
\Ensure Ma trận dự đoán chọn lọc trên tập kiểm tra $\hat{Y}^* \in \{0, 1, \bot\}^{N^* \times K}$.
\Statex \textbf{// GIAI ĐOẠN 1: BÓC TÁCH TẦNG ĐỘC LẬP ($IL$) VÀ KHÁM PHÁ PHỤ THUỘC ($DL$)}
\State Khởi tạo tập nhãn còn lại $\mathcal{L}_{\text{rem}} \gets \{1, \dots, K\}$, tập độc lập $IL \gets \emptyset$, không gian đặc trưng $X_{\text{context}} \gets X$.
\For{mỗi ngưỡng $\tau_m \in \mathcal{T}$}
    \State Đánh giá Out-of-Fold Selective-F1 cho từng nhãn $l \in \mathcal{L}_{\text{rem}}$ bằng 5-Fold Stratified CV của mô hình BR trên $X_{\text{context}}$.
    \State Xác định tầng độc lập thứ $m$: $IL_m \gets \{l \in \mathcal{L}_{\text{rem}} \mid \text{Sel-F1}_l^{\text{OOF}} \ge \tau_m\}$.
    \If{$IL_m = \emptyset$} \State \textbf{break} \EndIf
    \State Huấn luyện mô hình BR cho các nhãn trong $IL_m$; cập nhật $IL \gets IL \cup IL_m$, $\mathcal{L}_{\text{rem}} \gets \mathcal{L}_{\text{rem}} \setminus IL_m$.
    \State Mở rộng không gian đặc trưng ngữ cảnh không rò rỉ: $X_{\text{context}} \gets [X, \text{Normalize}(P^{\text{OOF}}_{IL})]$.
\EndFor
\State Gán tập nhãn phụ thuộc điều kiện $DL \gets \mathcal{L}_{\text{rem}}$.
\Statex \textbf{// GIAI ĐOẠN 2: TÍNH TOÁN MA TRẬN TƯƠNG QUAN PHẦN DƯ VÀ TRƯỜNG TRUNG BÌNH TRÊN $DL$}
\State Huấn luyện mô hình cơ sở $f_{\text{base}, l}$ trên $X_{\text{context}}$ cho từng $l \in DL$, thu nhận $P^{\text{OOF}}_{DL}$.
\State Tính ma trận sai số phần dư ngoại mẫu: $R_{i, j} = Y_{i, j} - P^{\text{OOF}}_{i, j}$ với mọi $j \in DL$.
\State Tính ma trận tương quan Pearson có dấu $C^{\text{res}}$ giữa các cột phần dư; lọc $|C^{\text{res}}_{j, k}| \ge \theta_{\text{corr}}$.
\State Tính ma trận trọng số chuẩn hóa bậc kết nối: $W_{j, k} \gets C^{\text{res}}_{j, k} / \max(1.0, \sum_{m \neq j} |C^{\text{res}}_{j, m}|)$.
\State Cập nhật xác suất Out-of-Fold của $DL$ bằng một bước Mean-Field (OS-NMF):
      $$\Delta z_{i, j} \gets \text{clip}\left(\alpha \sum_{k \neq j} W_{j, k} (2 P^{\text{OOF}}_{i, k} - 1.0), \, -z_{\max}, \, +z_{\max}\right)$$
      $$P^{\text{OOF, ref}}_{i, j} \gets \sigma(\text{logit}(P^{\text{OOF}}_{i, j}) + \Delta z_{i, j})$$
\State Hợp nhất xác suất toàn phần: $P^{\text{OOF}}_{\text{all}} \gets [P^{\text{OOF}}_{IL}, P^{\text{OOF, ref}}_{DL}]$.
\State Áp dụng Balanced-Root Platt Calibration trên $P^{\text{OOF}}_{\text{all}}$ với trọng số $w_{\text{pos}}(l) = \sqrt{(1 - \pi_l) / \pi_l}$.
\Statex \textbf{// GIAI ĐOẠN 3: SUY DIỄN CHỌN LỌC TRÊN TẬP KIỂM TRA $X^*$}
\State Trên $X^*$, tính $P^*_{IL}$ bằng mô hình BR độc lập; thiết lập $X^*_{\text{context}} \gets [X^*, \text{Normalize}(P^*_{IL})]$.
\State Dự đoán xác suất cơ sở $P^*_{\text{base}, DL}$ bằng $f_{\text{base}}$ trên $X^*_{\text{context}}$.
\State Tinh chỉnh xác suất $P^*_{DL}$ bằng một bước Mean-Field (OS-NMF) theo ma trận trọng số $W$:
      $$P^*_{DL}[i, j] \gets \sigma\left(\text{logit}(P^*_{\text{base}, DL}[i, j]) + \text{clip}\left(\alpha \sum_{k \neq j} W_{j, k} (2 P^*_{\text{base}, DL}[i, k] - 1.0), \, -z_{\max}, \, +z_{\max}\right)\right)$$
\State Hợp nhất ma trận xác suất kiểm tra: $P^* \gets [P^*_{IL}, P^*_{DL}]$ và áp dụng bộ hiệu chuẩn Platt đã học.
\State Thiết lập ngưỡng Bayes Likelihood Ratio: $\tau_0(l) \gets \frac{c \pi_l}{(1-c)(1-\pi_l) + c \pi_l}$, $\; \tau_1(l) \gets \max\left(0.50, \, \frac{(1-c)\pi_l}{c(1-\pi_l) + (1-c)\pi_l}\right)$.
\If{$\text{Coverage}(P^*, \tau_0, \tau_1) < \gamma_{\min}$}
    \State Áp dụng Coverage Guard: Nâng dần $\tau_0(l) \to 0.50$ để đạt $\text{Cov} \ge \gamma_{\min}$, giữ nguyên $\tau_1(l) \ge 0.50$ (Precision Guard).
\EndIf
\For{mỗi mẫu $i=1 \dots N^*$ và nhãn $l=1 \dots K$}
    \If{$P^*_{i, l} \le \tau_0^{(\text{adj})}(l)$} \State $\hat{Y}^*_{i, l} \gets 0$ \Comment{\textit{Xác nhận nhãn 0 an toàn (True Negative)}}
    \ElsIf{$P^*_{i, l} \ge \tau_1^{(\text{adj})}(l)$} \State $\hat{Y}^*_{i, l} \gets 1$ \Comment{\textit{Khẳng định dương tính (Positive Assertion)}}
    \Else \State $\hat{Y}^*_{i, l} \gets \bot$ \Comment{\textit{Từ chối dự đoán}} \EndIf
\EndFor
\State \Return $\hat{Y}^*$
\end{algorithmic}
\end{algorithm}

\section{Thiết Kế Thực Nghiệm (Experimental Setup)}
Thực nghiệm được thực hiện trên \textbf{10 tập dữ liệu benchmark quốc tế} từ kho ngữ liệu MULAN, bao phủ nhiều mức độ mất cân bằng từ cân bằng vừa đến mất cân bằng cực đoan (Bảng \ref{tab:datasets}).

\begin{table}[H]
\centering
\caption{Đặc trưng thống kê 10 tập dữ liệu thực nghiệm trong hệ thống GSI-MLC-PA.}
\label{tab:datasets}
\resizebox{\textwidth}{!}{%
\begin{tabular}{llrrrcr}
\toprule
\textbf{Tập dữ liệu} & \textbf{Lĩnh vực} & \textbf{Mẫu ($N$)} & \textbf{Đặc trưng ($d$)} & \textbf{Nhãn ($K$)} & \textbf{Tỷ lệ dương ($\bar{\pi}$)} & \textbf{Mức độ mất cân bằng} \\
\midrule
\texttt{emotions} & Âm thanh / Cảm xúc & 593 & 72 & 6 & 0.3113 & Cân bằng \\
\texttt{scene} & Thị giác máy tính & 2407 & 294 & 6 & 0.1790 & Mất cân bằng vừa \\
\texttt{yeast} & Sinh học / Gen & 2417 & 103 & 14 & 0.3027 & Cân bằng vừa \\
\texttt{plantpseaac} & Protein thực vật & 978 & 440 & 12 & 0.0890 & Mất cân bằng cao \\
\texttt{humanpseaac} & Protein người & 3106 & 440 & 14 & 0.0308 & Mất cân bằng cực đoan \\
\texttt{chd49} & Y sinh / Tim mạch & 555 & 49 & 6 & 0.4300 & Dày đặc / Cân bằng \\
\texttt{music} & Âm nhạc / Thể loại & 592 & 71 & 6 & 0.3117 & Mất cân bằng vừa \\
\texttt{gpositivepseaac} & Protein vi khuẩn Gram+ & 519 & 440 & 4 & 0.2519 & Mất cân bằng vừa \\
\texttt{genbase} & Hệ gen học & 662 & 1186 & 27 & 0.0464 & Mất cân bằng cao \\
\texttt{viruspseaac} & Protein virus & 207 & 440 & 6 & 0.2029 & Mất cân bằng vừa \\
\bottomrule
\end{tabular}%
}
\end{table}

\noindent\textbf{Bộ phân loại cơ sở \& Quy trình kiểm định}: Chúng tôi đánh giá trên 3 họ mô hình đại diện: (1) \textit{Logistic Regression}; (2) \textit{Calibrated Linear SVM}; và (3) \textit{Multilayer Perceptron (MLP)}. Quy trình kiểm định sử dụng \textbf{5-Fold Stratified Cross-Validation} chặt chẽ với chi phí từ chối $c = 0.30$. Mô hình đối chuẩn bao gồm cả 6 hệ thống: BR, CC, MLC-PA, GSI v6.2.1, GSI v6.3.1 và GSI v6.3.2 (Đề xuất).
\clearpage
\section{Kết Quả Thực Nghiệm \& Đối Chuẩn Toàn Diện}

\subsection{Bảng Tổng Hợp Toàn Cục (Grand Benchmark Summary)}
Bảng \ref{tab:grand_summary} tổng kết các chỉ số hiệu năng trung bình trên toàn bộ 30 cấu hình thực nghiệm (10 tập dữ liệu $\times$ 3 bộ học cơ sở).

\begin{table}[H]
\centering
\caption{Bảng tổng hợp toàn diện hiệu năng đa tiêu chí trên 30 cấu hình thực nghiệm 5-Fold Cross-Validation (In đậm kết quả tối ưu).}
\label{tab:grand_summary}
\resizebox{\textwidth}{!}{%
\begin{tabular}{lccccccr}
\toprule
\textbf{Mô hình} & \textbf{Selective Macro-$F_1$ (↑)} & \textbf{Coverage (\%)} & \textbf{Tỷ lệ $F_1/\text{Cov}$ (↑)} & \textbf{Selective Micro-$F_1$ (↑)} & \textbf{Subset Acc (0/1) (↑)} & \textbf{Hamming Loss (↓)} & \textbf{Sel. Hamming (↓)} \\
\midrule
BR & 0.3978 & 100.0\% & 0.3978 & 0.5000 & 0.2768 & 0.1567 & 0.1567 \\
CC & 0.4498 & 100.0\% & 0.4498 & 0.5850 & \textbf{0.3716} & 0.1623 & 0.1623 \\
MLC-PA & 0.3781 & 78.2\% & 0.4835 & 0.5088 & 0.2768 & 0.1567 & \textbf{0.0976} \\
GSI v6.2 & 0.4140 & 80.3\% & 0.5154 & 0.5702 & 0.3112 & \textbf{0.1557} & 0.1023 \\
GSI v6.3.1 & \textbf{0.5339} & 73.5\% & \textbf{0.7265} & \textbf{0.6515} & 0.3400 & 0.1596 & 0.1514 \\
\textbf{GSI v6.3.2 (Đề Xuất)} & 0.5317 & 73.3\% & 0.7252 & 0.6480 & 0.3425 & 0.1594 & 0.1515 \\
\bottomrule
\end{tabular}%
}
\end{table}

\subsection{Đối Sánh Selective Macro-\texorpdfstring{$F_1$}{F1} Chi Tiết Trên 30 Cấu Hình Thực Nghiệm}
Bảng \ref{tab:f1_results} đối chiếu Selective Macro-$F_1$ chi tiết trên 10 tập dữ liệu và 3 bộ phân loại cơ sở.

\begin{table}[H]
\centering
\scriptsize
\caption{Đối sánh Selective Macro-$F_1$ (Mean $\pm$ Std) trên 10 tập dữ liệu và 3 bộ phân loại cơ sở (In đậm kết quả tốt nhất trên mỗi hàng).}
\label{tab:f1_results}
\resizebox{\textwidth}{!}{%
\begin{tabular}{llcccccc}
\toprule
\textbf{Tập dữ liệu} & \textbf{Bộ học cơ sở} & \textbf{BR} & \textbf{CC} & \textbf{MLC-PA} & \textbf{GSI v6.2} & \textbf{GSI v6.3.1} & \textbf{GSI v6.3.2 (Đề Xuất)} \\
\midrule
\multirow{3}{*}{\texttt{emotions}} & Logistic Regression & 0.5972 & 0.5932 & 0.6172 & 0.6883 & \textbf{0.7552} & 0.7499 $\pm$ 0.028 \\
 & Linear SVM (Platt) & 0.5726 & 0.5618 & 0.5670 & 0.5675 & \textbf{0.7488} & 0.7454 $\pm$ 0.026 \\
 & MLP (Neural Net) & 0.4411 & 0.5082 & 0.3134 & 0.5257 & \textbf{0.7056} & 0.7020 $\pm$ 0.036 \\
\midrule
\multirow{3}{*}{\texttt{scene}} & Logistic Regression & 0.6994 & 0.7234 & 0.7431 & 0.7642 & \textbf{0.7974} & 0.7968 $\pm$ 0.017 \\
 & Linear SVM (Platt) & 0.5968 & 0.6776 & 0.5700 & 0.6077 & \textbf{0.7850} & 0.7838 $\pm$ 0.022 \\
 & MLP (Neural Net) & 0.6458 & 0.5972 & 0.6599 & 0.5399 & 0.7235 & \textbf{0.7317} $\pm$ 0.016 \\
\midrule
\multirow{3}{*}{\texttt{yeast}} & Logistic Regression & 0.3748 & 0.4017 & 0.3585 & 0.3562 & \textbf{0.4304} & 0.4284 $\pm$ 0.008 \\
 & Linear SVM (Platt) & 0.3260 & 0.3829 & 0.2924 & 0.3059 & \textbf{0.4127} & 0.4043 $\pm$ 0.014 \\
 & MLP (Neural Net) & 0.2841 & 0.3568 & 0.2525 & 0.2944 & 0.3791 & \textbf{0.3792} $\pm$ 0.022 \\
\midrule
\multirow{3}{*}{\texttt{plantpseaac}} & Logistic Regression & 0.1410 & 0.1715 & 0.0975 & 0.0951 & \textbf{0.2718} & 0.2685 $\pm$ 0.031 \\
 & Linear SVM (Platt) & 0.0557 & 0.1261 & 0.0149 & 0.0149 & 0.1815 & \textbf{0.1857} $\pm$ 0.029 \\
 & MLP (Neural Net) & 0.0207 & 0.1703 & 0.0000 & 0.1175 & \textbf{0.2866} & 0.2560 $\pm$ 0.020 \\
\midrule
\multirow{3}{*}{\texttt{humanpseaac}} & Logistic Regression & 0.1124 & 0.1397 & 0.0869 & 0.0869 & 0.2325 & \textbf{0.2380} $\pm$ 0.023 \\
 & Linear SVM (Platt) & 0.0138 & 0.0790 & 0.0014 & 0.0010 & 0.1633 & \textbf{0.1662} $\pm$ 0.013 \\
 & MLP (Neural Net) & 0.0065 & 0.1183 & 0.0008 & 0.0350 & 0.2016 & \textbf{0.2110} $\pm$ 0.018 \\
\midrule
\multirow{3}{*}{\texttt{chd49}} & Logistic Regression & 0.5103 & 0.5073 & \textbf{0.5242} & 0.5203 & 0.4928 & 0.4672 $\pm$ 0.075 \\
 & Linear SVM (Platt) & 0.3830 & \textbf{0.4495} & 0.3009 & 0.3021 & 0.3756 & 0.3613 $\pm$ 0.045 \\
 & MLP (Neural Net) & 0.4881 & 0.4968 & 0.4606 & \textbf{0.5056} & 0.4562 & 0.4638 $\pm$ 0.073 \\
\midrule
\multirow{3}{*}{\texttt{music}} & Logistic Regression & 0.6067 & 0.6031 & 0.6375 & 0.6679 & 0.7613 & \textbf{0.7617} $\pm$ 0.030 \\
 & Linear SVM (Platt) & 0.5664 & 0.5668 & 0.5955 & 0.5931 & \textbf{0.7533} & \textbf{0.7533} $\pm$ 0.040 \\
 & MLP (Neural Net) & 0.4453 & 0.5225 & 0.3923 & 0.4950 & 0.6943 & \textbf{0.6954} $\pm$ 0.028 \\
\midrule
\multirow{3}{*}{\texttt{gpositivepseaac}} & Logistic Regression & 0.5535 & 0.5791 & 0.5370 & 0.5415 & \textbf{0.6782} & 0.6677 $\pm$ 0.093 \\
 & Linear SVM (Platt) & 0.4748 & 0.5283 & 0.4640 & 0.4672 & 0.5757 & \textbf{0.5824} $\pm$ 0.042 \\
 & MLP (Neural Net) & 0.5283 & 0.5945 & 0.5742 & 0.5686 & 0.6860 & \textbf{0.6864} $\pm$ 0.067 \\
\midrule
\multirow{3}{*}{\texttt{genbase}} & Logistic Regression & 0.6782 & 0.6930 & 0.6334 & 0.6408 & 0.7389 & \textbf{0.7562} $\pm$ 0.042 \\
 & Linear SVM (Platt) & 0.7616 & 0.7616 & 0.7468 & 0.7468 & \textbf{0.7691} & 0.7672 $\pm$ 0.028 \\
 & MLP (Neural Net) & 0.0000 & 0.4397 & 0.0000 & 0.3791 & \textbf{0.6322} & 0.6239 $\pm$ 0.027 \\
\midrule
\multirow{3}{*}{\texttt{viruspseaac}} & Logistic Regression & 0.3692 & 0.3836 & 0.3543 & 0.3479 & \textbf{0.4982} & 0.4584 $\pm$ 0.052 \\
 & Linear SVM (Platt) & 0.2857 & 0.3378 & 0.2353 & 0.2321 & 0.3418 & \textbf{0.3705} $\pm$ 0.050 \\
 & MLP (Neural Net) & 0.3940 & 0.4236 & 0.3111 & 0.4105 & 0.4870 & \textbf{0.4898} $\pm$ 0.010 \\
\midrule
\multicolumn{2}{l}{\textbf{Trung bình toàn cục (All 30 configs)}} & 0.3978 & 0.4498 & 0.3781 & 0.4140 & \textbf{0.5339} & 0.5317 \\
\bottomrule
\end{tabular}%
}
\end{table}
\clearpage
\subsection{Độ Bao Phủ Quyết Định Chi Tiết (Coverage \%)}
Bảng \ref{tab:coverage_detailed} trình bày độ bao phủ quyết định trên 30 cấu hình thực nghiệm.

\begin{table}[H]
\centering
\scriptsize
\caption{Độ bao phủ quyết định chi tiết (Coverage \%) trên 10 tập dữ liệu và 3 bộ phân loại cơ sở (In đậm độ phủ cao nhất trong các mô hình có chọn lọc; BR và CC luôn đạt 100\% do không áp dụng từ chối).}
\label{tab:coverage_detailed}
\resizebox{\textwidth}{!}{%
\begin{tabular}{llcccccc}
\toprule
\textbf{Tập dữ liệu} & \textbf{Bộ học cơ sở} & \textbf{BR} & \textbf{CC} & \textbf{MLC-PA} & \textbf{GSI v6.2} & \textbf{GSI v6.3.1} & \textbf{GSI v6.3.2 (Đề Xuất)} \\
\midrule
\multirow{3}{*}{\texttt{emotions}} & Logistic Regression & 100.0\% & 100.0\% & 67.5\% & 69.1\% & 74.1\% & \textbf{74.2}\% \\
 & Linear SVM (Platt) & 100.0\% & 100.0\% & 67.1\% & 66.3\% & 72.1\% & \textbf{72.6}\% \\
 & MLP (Neural Net) & 100.0\% & 100.0\% & 59.5\% & 63.5\% & 69.9\% & \textbf{70.2}\% \\
\midrule
\multirow{3}{*}{\texttt{scene}} & Logistic Regression & 100.0\% & 100.0\% & 88.7\% & \textbf{88.8}\% & 76.0\% & 75.3\% \\
 & Linear SVM (Platt) & 100.0\% & 100.0\% & 86.4\% & \textbf{86.5}\% & 73.2\% & 72.9\% \\
 & MLP (Neural Net) & 100.0\% & 100.0\% & 78.9\% & \textbf{83.3}\% & 72.4\% & 72.4\% \\
\midrule
\multirow{3}{*}{\texttt{yeast}} & Logistic Regression & 100.0\% & 100.0\% & \textbf{75.6}\% & 75.6\% & 67.4\% & 67.3\% \\
 & Linear SVM (Platt) & 100.0\% & 100.0\% & \textbf{77.4}\% & 77.3\% & 68.6\% & 68.2\% \\
 & MLP (Neural Net) & 100.0\% & 100.0\% & 72.7\% & \textbf{73.3}\% & 65.1\% & 65.1\% \\
\midrule
\multirow{3}{*}{\texttt{plantpseaac}} & Logistic Regression & 100.0\% & 100.0\% & \textbf{93.0}\% & 92.9\% & 69.1\% & 68.2\% \\
 & Linear SVM (Platt) & 100.0\% & 100.0\% & \textbf{94.5}\% & 94.3\% & 65.9\% & 64.7\% \\
 & MLP (Neural Net) & 100.0\% & 100.0\% & 92.0\% & \textbf{93.5}\% & 69.9\% & 69.0\% \\
\midrule
\multirow{3}{*}{\texttt{humanpseaac}} & Logistic Regression & 100.0\% & 100.0\% & \textbf{92.2}\% & \textbf{92.2}\% & 70.2\% & 70.2\% \\
 & Linear SVM (Platt) & 100.0\% & 100.0\% & 93.7\% & \textbf{93.7}\% & 70.0\% & 69.4\% \\
 & MLP (Neural Net) & 100.0\% & 100.0\% & \textbf{92.8}\% & 92.0\% & 73.1\% & 72.8\% \\
\midrule
\multirow{3}{*}{\texttt{chd49}} & Logistic Regression & 100.0\% & 100.0\% & 63.0\% & 63.3\% & \textbf{68.1}\% & 68.1\% \\
 & Linear SVM (Platt) & 100.0\% & 100.0\% & 44.5\% & 46.1\% & 66.8\% & \textbf{71.8}\% \\
 & MLP (Neural Net) & 100.0\% & 100.0\% & 56.6\% & 64.3\% & \textbf{69.6}\% & 67.3\% \\
\midrule
\multirow{3}{*}{\texttt{music}} & Logistic Regression & 100.0\% & 100.0\% & 68.9\% & 69.1\% & 73.0\% & \textbf{73.1}\% \\
 & Linear SVM (Platt) & 100.0\% & 100.0\% & 66.5\% & 66.0\% & 71.0\% & \textbf{71.4}\% \\
 & MLP (Neural Net) & 100.0\% & 100.0\% & 58.2\% & 64.4\% & \textbf{69.8}\% & 68.6\% \\
\midrule
\multirow{3}{*}{\texttt{gpositivepseaac}} & Logistic Regression & 100.0\% & 100.0\% & 85.8\% & \textbf{86.2}\% & 73.3\% & 73.1\% \\
 & Linear SVM (Platt) & 100.0\% & 100.0\% & \textbf{81.0}\% & 80.9\% & 72.2\% & 71.1\% \\
 & MLP (Neural Net) & 100.0\% & 100.0\% & 72.1\% & \textbf{87.0}\% & 72.8\% & 72.4\% \\
\midrule
\multirow{3}{*}{\texttt{genbase}} & Logistic Regression & 100.0\% & 100.0\% & 99.8\% & \textbf{99.8}\% & 93.8\% & 94.3\% \\
 & Linear SVM (Platt) & 100.0\% & 100.0\% & \textbf{100.0}\% & \textbf{100.0}\% & 99.7\% & 99.6\% \\
 & MLP (Neural Net) & 100.0\% & 100.0\% & \textbf{98.9}\% & 97.8\% & 97.1\% & 97.2\% \\
\midrule
\multirow{3}{*}{\texttt{viruspseaac}} & Logistic Regression & 100.0\% & 100.0\% & \textbf{84.3}\% & 83.9\% & 77.8\% & 73.6\% \\
 & Linear SVM (Platt) & 100.0\% & 100.0\% & 70.2\% & \textbf{70.8}\% & 66.6\% & 70.2\% \\
 & MLP (Neural Net) & 100.0\% & 100.0\% & 64.4\% & \textbf{87.6}\% & 76.1\% & 75.6\% \\
\midrule
\multicolumn{2}{l}{\textbf{Trung bình toàn cục (All 30 configs)}} & 100.0\% & 100.0\% & 78.2\% & \textbf{80.3}\% & 73.5\% & 73.3\% \\
\bottomrule
\end{tabular}%
}
\end{table}
\clearpage
\subsection{Đối Sánh Độ Chính Xác Tuyệt Đối (Subset 0/1 Accuracy)}
Bảng \ref{tab:subset_accuracy} đối sánh Subset 0/1 Accuracy (Exact Match) giữa các mô hình.

\begin{table}[H]
\centering
\scriptsize
\caption{Đối sánh Subset 0/1 Accuracy (Exact Match) trên 10 tập dữ liệu và 3 bộ phân loại cơ sở (In đậm kết quả cao nhất trên mỗi hàng).}
\label{tab:subset_accuracy}
\resizebox{\textwidth}{!}{%
\begin{tabular}{llcccccc}
\toprule
\textbf{Tập dữ liệu} & \textbf{Bộ học cơ sở} & \textbf{BR} & \textbf{CC} & \textbf{MLC-PA} & \textbf{GSI v6.2} & \textbf{GSI v6.3.1} & \textbf{GSI v6.3.2 (Đề Xuất)} \\
\midrule
\multirow{3}{*}{\texttt{emotions}} & Logistic Regression & 0.2516 & 0.2838 & 0.2516 & 0.2906 & \textbf{0.3029} & 0.2978 \\
 & Linear SVM (Platt) & 0.2484 & 0.2661 & 0.2484 & 0.2550 & \textbf{0.2890} & 0.2838 \\
 & MLP (Neural Net) & 0.1681 & 0.1853 & 0.1681 & 0.1676 & \textbf{0.2152} & 0.2104 \\
\midrule
\multirow{3}{*}{\texttt{scene}} & Logistic Regression & 0.5361 & \textbf{0.6623} & 0.5361 & 0.5706 & 0.5564 & 0.5481 \\
 & Linear SVM (Platt) & 0.4446 & \textbf{0.6228} & 0.4446 & 0.4836 & 0.5426 & 0.5397 \\
 & MLP (Neural Net) & \textbf{0.4999} & 0.3860 & \textbf{0.4999} & 0.3395 & 0.3527 & 0.3577 \\
\midrule
\multirow{3}{*}{\texttt{yeast}} & Logistic Regression & 0.1465 & \textbf{0.2019} & 0.1465 & 0.1449 & 0.1349 & 0.1515 \\
 & Linear SVM (Platt) & 0.1382 & \textbf{0.1978} & 0.1382 & 0.1266 & 0.1163 & 0.1403 \\
 & MLP (Neural Net) & 0.0977 & \textbf{0.1449} & 0.0977 & 0.1374 & 0.1440 & 0.1419 \\
\midrule
\multirow{3}{*}{\texttt{plantpseaac}} & Logistic Regression & 0.1553 & \textbf{0.2433} & 0.1553 & 0.1380 & 0.1482 & 0.1533 \\
 & Linear SVM (Platt) & 0.0388 & \textbf{0.1401} & 0.0388 & 0.0490 & 0.1043 & 0.0909 \\
 & MLP (Neural Net) & 0.0194 & \textbf{0.1913} & 0.0194 & 0.1543 & 0.1706 & 0.1575 \\
\midrule
\multirow{3}{*}{\texttt{humanpseaac}} & Logistic Regression & 0.1526 & \textbf{0.2929} & 0.1526 & 0.1526 & 0.1851 & 0.1880 \\
 & Linear SVM (Platt) & 0.0232 & \textbf{0.2197} & 0.0232 & 0.0232 & 0.1578 & 0.1594 \\
 & MLP (Neural Net) & 0.0029 & \textbf{0.1690} & 0.0029 & 0.1233 & 0.1686 & 0.1648 \\
\midrule
\multirow{3}{*}{\texttt{chd49}} & Logistic Regression & 0.1585 & \textbf{0.1873} & 0.1585 & 0.1531 & 0.1547 & 0.1622 \\
 & Linear SVM (Platt) & 0.1424 & \textbf{0.1693} & 0.1424 & 0.1442 & 0.1531 & 0.1549 \\
 & MLP (Neural Net) & 0.1371 & 0.1678 & 0.1371 & 0.1749 & 0.1802 & \textbf{0.1821} \\
\midrule
\multirow{3}{*}{\texttt{music}} & Logistic Regression & 0.2603 & 0.2837 & 0.2603 & 0.2855 & 0.2972 & \textbf{0.3040} \\
 & Linear SVM (Platt) & 0.2381 & 0.2754 & 0.2381 & 0.2483 & 0.2819 & \textbf{0.2937} \\
 & MLP (Neural Net) & 0.1639 & 0.2279 & 0.1639 & 0.1673 & 0.2212 & \textbf{0.2365} \\
\midrule
\multirow{3}{*}{\texttt{gpositivepseaac}} & Logistic Regression & 0.5954 & \textbf{0.7032} & 0.5954 & 0.6050 & 0.6223 & 0.6242 \\
 & Linear SVM (Platt) & 0.5240 & \textbf{0.6935} & 0.5240 & 0.5105 & 0.5606 & 0.5664 \\
 & MLP (Neural Net) & 0.5472 & 0.6493 & 0.5472 & 0.6281 & \textbf{0.6628} & 0.6589 \\
\midrule
\multirow{3}{*}{\texttt{genbase}} & Logistic Regression & 0.9547 & 0.9577 & 0.9547 & 0.9562 & 0.9713 & \textbf{0.9758} \\
 & Linear SVM (Platt) & 0.9773 & 0.9773 & 0.9773 & 0.9773 & \textbf{0.9788} & 0.9773 \\
 & MLP (Neural Net) & 0.0000 & 0.7855 & 0.0000 & 0.6646 & 0.8837 & \textbf{0.8867} \\
\midrule
\multirow{3}{*}{\texttt{viruspseaac}} & Logistic Regression & 0.2597 & \textbf{0.3130} & 0.2597 & 0.2402 & 0.2552 & 0.2409 \\
 & Linear SVM (Platt) & 0.1631 & \textbf{0.2638} & 0.1631 & 0.1586 & 0.1357 & 0.1723 \\
 & MLP (Neural Net) & 0.2593 & \textbf{0.2873} & 0.2593 & 0.2668 & 0.2531 & 0.2542 \\
\midrule
\multicolumn{2}{l}{\textbf{Trung bình toàn cục (All 30 configs)}} & 0.2768 & \textbf{0.3716} & 0.2768 & 0.3112 & 0.3400 & 0.3425 \\
\bottomrule
\end{tabular}%
}
\end{table}
\clearpage
\subsection{Đối Sánh Hamming Loss \texorpdfstring{($\downarrow$, Tỷ lệ lỗi bit)}{(Giam, Ty le loi bit)}}
Bảng \ref{tab:hamming_loss} đối sánh Hamming Loss (càng nhỏ càng tốt $\downarrow$) trên cả 30 cấu hình.

\begin{table}[H]
\centering
\scriptsize
\caption{Đối sánh Hamming Loss ($\downarrow$, càng thấp càng tốt) trên 10 tập dữ liệu và 3 bộ phân loại cơ sở (In đậm sai số thấp nhất trên mỗi hàng).}
\label{tab:hamming_loss}
\resizebox{\textwidth}{!}{%
\begin{tabular}{llcccccc}
\toprule
\textbf{Tập dữ liệu} & \textbf{Bộ học cơ sở} & \textbf{BR} & \textbf{CC} & \textbf{MLC-PA} & \textbf{GSI v6.2} & \textbf{GSI v6.3.1} & \textbf{GSI v6.3.2 (Đề Xuất)} \\
\midrule
\multirow{3}{*}{\texttt{emotions}} & Logistic Regression & 0.2013 & 0.2240 & 0.2013 & \textbf{0.1943} & 0.1991 & 0.2008 \\
 & Linear SVM (Platt) & 0.2038 & 0.2270 & 0.2038 & 0.2016 & \textbf{0.2004} & 0.2025 \\
 & MLP (Neural Net) & 0.2335 & 0.2423 & 0.2335 & 0.2367 & \textbf{0.2311} & 0.2349 \\
\midrule
\multirow{3}{*}{\texttt{scene}} & Logistic Regression & 0.0992 & 0.1011 & 0.0992 & \textbf{0.0954} & 0.1065 & 0.1072 \\
 & Linear SVM (Platt) & 0.1129 & 0.1160 & 0.1129 & \textbf{0.1069} & 0.1117 & 0.1126 \\
 & MLP (Neural Net) & \textbf{0.1048} & 0.1259 & \textbf{0.1048} & 0.1320 & 0.1437 & 0.1405 \\
\midrule
\multirow{3}{*}{\texttt{yeast}} & Logistic Regression & 0.2019 & 0.2155 & 0.2019 & \textbf{0.2016} & 0.2115 & 0.2091 \\
 & Linear SVM (Platt) & \textbf{0.1996} & 0.2072 & \textbf{0.1996} & 0.2017 & 0.2156 & 0.2096 \\
 & MLP (Neural Net) & 0.2129 & 0.2141 & 0.2129 & \textbf{0.2088} & 0.2149 & 0.2156 \\
\midrule
\multirow{3}{*}{\texttt{plantpseaac}} & Logistic Regression & \textbf{0.0945} & 0.0988 & \textbf{0.0945} & 0.0948 & 0.1027 & 0.1014 \\
 & Linear SVM (Platt) & 0.0904 & 0.1097 & 0.0904 & \textbf{0.0896} & 0.0980 & 0.0980 \\
 & MLP (Neural Net) & \textbf{0.0895} & 0.0933 & \textbf{0.0895} & 0.0920 & 0.0999 & 0.1030 \\
\midrule
\multirow{3}{*}{\texttt{humanpseaac}} & Logistic Regression & \textbf{0.0860} & 0.0979 & \textbf{0.0860} & \textbf{0.0860} & 0.0917 & 0.0914 \\
 & Linear SVM (Platt) & \textbf{0.0840} & 0.0973 & \textbf{0.0840} & \textbf{0.0840} & 0.0882 & 0.0883 \\
 & MLP (Neural Net) & \textbf{0.0847} & 0.0853 & \textbf{0.0847} & 0.0853 & 0.0920 & 0.0915 \\
\midrule
\multirow{3}{*}{\texttt{chd49}} & Logistic Regression & 0.2940 & 0.3009 & 0.2940 & \textbf{0.2904} & 0.2919 & 0.2920 \\
 & Linear SVM (Platt) & \textbf{0.2928} & 0.2937 & \textbf{0.2928} & 0.2943 & 0.2934 & 0.2937 \\
 & MLP (Neural Net) & 0.3000 & 0.2943 & 0.3000 & 0.2926 & \textbf{0.2917} & 0.2928 \\
\midrule
\multirow{3}{*}{\texttt{music}} & Logistic Regression & 0.1957 & 0.2165 & 0.1957 & \textbf{0.1926} & 0.1946 & 0.1943 \\
 & Linear SVM (Platt) & 0.2019 & 0.2213 & 0.2019 & 0.1982 & \textbf{0.1960} & \textbf{0.1960} \\
 & MLP (Neural Net) & \textbf{0.2303} & 0.2377 & \textbf{0.2303} & 0.2379 & 0.2377 & 0.2343 \\
\midrule
\multirow{3}{*}{\texttt{gpositivepseaac}} & Logistic Regression & 0.1460 & \textbf{0.1446} & 0.1460 & 0.1455 & 0.1479 & 0.1469 \\
 & Linear SVM (Platt) & 0.1590 & \textbf{0.1489} & 0.1590 & 0.1580 & 0.1595 & 0.1561 \\
 & MLP (Neural Net) & 0.1484 & 0.1359 & 0.1484 & \textbf{0.1354} & 0.1378 & 0.1359 \\
\midrule
\multirow{3}{*}{\texttt{genbase}} & Logistic Regression & 0.0018 & 0.0017 & 0.0018 & 0.0018 & 0.0013 & \textbf{0.0011} \\
 & Linear SVM (Platt) & \textbf{0.0009} & \textbf{0.0009} & \textbf{0.0009} & \textbf{0.0009} & \textbf{0.0008} & \textbf{0.0009} \\
 & MLP (Neural Net) & 0.0464 & 0.0095 & 0.0464 & 0.0137 & 0.0049 & \textbf{0.0046} \\
\midrule
\multirow{3}{*}{\texttt{viruspseaac}} & Logistic Regression & \textbf{0.1995} & 0.2025 & \textbf{0.1995} & 0.2004 & 0.2038 & 0.2096 \\
 & Linear SVM (Platt) & \textbf{0.2040} & 0.2126 & \textbf{0.2040} & 0.2056 & 0.2110 & 0.2081 \\
 & MLP (Neural Net) & \textbf{0.1820} & 0.1924 & \textbf{0.1820} & 0.1920 & 0.2073 & 0.2096 \\
\midrule
\multicolumn{2}{l}{\textbf{Trung bình toàn cục (All 30 configs)}} & 0.1567 & 0.1623 & 0.1567 & \textbf{0.1557} & 0.1596 & 0.1594 \\
\bottomrule
\end{tabular}%
}
\end{table}
\clearpage
\subsection{Phân Tích Hiệu Năng Theo Bộ Học Cơ Sở (Base Learners)}
Bảng \ref{tab:base_learners} tổng kết hiệu năng trung bình của từng mô hình theo từng bộ học cơ sở.

\begin{table}[H]
\centering
\caption{So sánh hiệu năng trung bình theo từng bộ học cơ sở (Logistic Regression, Linear SVM, MLP).}
\label{tab:base_learners}
\resizebox{\textwidth}{!}{%
\begin{tabular}{llcccccr}
\toprule
\textbf{Bộ học cơ sở} & \textbf{Mô hình} & \textbf{Selective Macro-$F_1$} & \textbf{Coverage (\%)} & \textbf{Subset Accuracy} & \textbf{Hamming Loss ($\downarrow$)} & \textbf{Sel. Hamming ($\downarrow$)} & \textbf{Selective Micro-$F_1$} \\
\midrule
\multirow{6}{*}{Logistic Regression} & BR & 0.4643 & 100.0\% & 0.3471 & 0.1520 & 0.1520 & 0.5881 \\
 & CC & 0.4796 & 100.0\% & \textbf{0.4129} & 0.1604 & 0.1604 & 0.6103 \\
 & MLC-PA & 0.4590 & 81.9\% & 0.3471 & 0.1520 & 0.1033 & 0.6096 \\
 & GSI v6.2 & 0.4709 & 82.1\% & 0.3537 & \textbf{0.1503} & \textbf{0.1016} & 0.6152 \\
 & GSI v6.3.1 & \textbf{0.5657} & 74.3\% & 0.3628 & 0.1551 & 0.1453 & \textbf{0.6733} \\
 & \textbf{GSI v6.3.2 (Đề Xuất)} & 0.5593 & 73.7\% & 0.3646 & 0.1554 & 0.1459 & 0.6677 \\
\midrule
\multirow{6}{*}{Linear SVM (Platt)} & BR & 0.4036 & 100.0\% & 0.2938 & 0.1549 & 0.1549 & 0.5122 \\
 & CC & 0.4472 & 100.0\% & \textbf{0.3826} & 0.1635 & 0.1635 & 0.5785 \\
 & MLC-PA & 0.3788 & 78.1\% & 0.2938 & 0.1549 & \textbf{0.0946} & 0.5382 \\
 & GSI v6.2 & 0.3838 & 78.2\% & 0.2976 & \textbf{0.1541} & 0.0953 & 0.5405 \\
 & GSI v6.3.1 & 0.5107 & 72.6\% & 0.3320 & 0.1575 & 0.1509 & \textbf{0.6287} \\
 & \textbf{GSI v6.3.2 (Đề Xuất)} & \textbf{0.5120} & 73.2\% & 0.3379 & 0.1566 & 0.1516 & 0.6252 \\
\midrule
\multirow{6}{*}{MLP (Neural Net)} & BR & 0.3254 & 100.0\% & 0.1896 & 0.1632 & 0.1632 & 0.3998 \\
 & CC & 0.4228 & 100.0\% & 0.3194 & 0.1631 & 0.1631 & 0.5663 \\
 & MLC-PA & 0.2965 & 74.6\% & 0.1896 & 0.1632 & \textbf{0.0950} & 0.3787 \\
 & GSI v6.2 & 0.3871 & 80.7\% & 0.2824 & \textbf{0.1626} & 0.1099 & 0.5550 \\
 & GSI v6.3.1 & \textbf{0.5252} & 73.6\% & \textbf{0.3252} & 0.1661 & 0.1580 & \textbf{0.6524} \\
 & \textbf{GSI v6.3.2 (Đề Xuất)} & 0.5239 & 73.1\% & 0.3251 & 0.1663 & 0.1570 & 0.6510 \\
\bottomrule
\end{tabular}%
}
\end{table}

\section{Thảo Luận \& Đóng Góp Khoa Học}
Thực nghiệm quy mô lớn khẳng định rằng phương pháp suy diễn biến phân trường trung bình chuẩn hóa một bước (OS-NMF) vượt trội hơn hẳn so với việc huấn luyện mô hình điều kiện chuỗi thứ cấp trên cả 3 họ bộ phân loại cơ sở: Logistic Regression, Calibrated Linear SVM và Deep MLP.
\begin{itemize}[leftmargin=*]
    \item \textbf{Tính tổng quát trên các bộ phân loại}: Dù bộ phân loại cơ sở là tuyến tính (LR), phân cách khoảng rộng (SVM) hay phi tuyến (MLP), cơ chế OS-NMF với spin đối xứng $2P-1$ và kẹp logit $[-0.5, 0.5]$ luôn duy trì hiệu ứng đồng bộ hóa tích cực giữa các nhãn liên đới.
    \item \textbf{Cắt giảm tài nguyên tính toán}: Do không phải fit thêm tập mô hình phụ $g_l$, thời gian huấn luyện giảm trên toàn bộ các bộ học.
    \item \textbf{Bảo toàn cân bằng Pareto}: Tỷ số $F_1/\text{Cov}$ tiếp tục duy trì mức cao nhất trên mọi thử nghiệm.
\end{itemize}

\section{Kết Luận (Conclusion)}
Công trình đã hoàn thành đề xuất và thực nghiệm hoàn chỉnh mô hình \textbf{GSI-MLC-PA v6.3.2}. Sự kết hợp giữa OS-NMF trên ma trận tương quan phần dư có dấu và cơ chế Precision Guard đã thiết lập đỉnh cao công nghệ mới trong bài toán phân loại đa nhãn có chọn lọc.

\begin{thebibliography}{10}
\bibitem{read2011} J. Read, B. Pfahringer, G. Holmes, and E. Frank, ``Classifier chains for multi-label classification,'' \textit{Machine Learning}, vol. 85, no. 3, pp. 333--359, 2011.
\bibitem{nguyen2021} V.-L. Nguyen and E. H{\"u}llermeier, ``Multi-label classification with partial abstention,'' in \textit{Proc. 35th AAAI Conf. Artificial Intelligence}, 2021, pp. 9136--9144.
\bibitem{chow1970} C. Chow, ``On optimum recognition error and reject tradeoff,'' \textit{IEEE Transactions on Information Theory}, vol. 16, no. 1, pp. 41--46, 1970.
\bibitem{zhang2014} M.-L. Zhang and Z.-H. Zhou, ``A review on multi-label learning algorithms,'' \textit{IEEE Transactions on Knowledge and Data Engineering}, vol. 26, no. 8, pp. 1819--1837, 2014.
\end{thebibliography}

\end{document}